\documentclass[11pt,a4paper]{article}
\usepackage[T1]{fontenc}
\usepackage{lmodern,amsmath,amssymb,mathtools,booktabs,tabularx,microtype,xcolor}
\usepackage[bookmarks=false]{hyperref}
\usepackage[margin=24mm,headheight=14pt]{geometry}
\usepackage{fancyhdr}
\definecolor{ink}{HTML}{244B68}
\hypersetup{colorlinks=true,linkcolor=ink,pdftitle={Dunnhumby: Expanding the Candidate Score}}
\pagestyle{fancy}\fancyhf{}\fancyhead[L]{\small Dunnhumby: expanding the score}\fancyfoot[C]{\thepage}
\setlength{\parindent}{0pt}\setlength{\parskip}{5pt}\setlength{\emergencystretch}{2em}
\newcommand{\I}{\mathcal I}\newcommand{\pos}[1]{\left[#1\right]_+}
\newcommand{\off}{\operatorname{off}}\newcommand{\diag}{\operatorname{diag}}
\newcommand{\N}{\mathcal N}\newcommand{\sd}{\operatorname{sd}}
\newcommand{\why}{\par\textbf{Purpose and intuition.} }
\numberwithin{equation}{section}
\begin{document}
\begin{center}
{\LARGE\textcolor{ink}{Dunnhumby: Expanding the Candidate Score}}\\[6pt]
{\large From $\theta(c\mid S)$ to training baskets, households, and product text}
\end{center}\medskip
Take a candidate product $c$ and a nonempty partially filled basket $S$, with $c\notin S$.
Let $\I$ be the known product catalogue, $n=|\I|$, and let
$B\in\{0,1\}^{m\times n}$ describe the training baskets: $B_{ei}=1$ when training
basket $e$ contains product $i$. Define $f_i=\sum_e B_{ei}$, its number of training-basket
appearances. Each basket also has a household label $u(e)$. Product text and a product
category label $\ell(i)$ are additional inputs, not quantities determined by $B$.

This note describes the Dunnhumby basket-completion formulation with the full revealed set
$S$. It starts from the score, expands its terms, and ends with the rules selecting its
fixed settings. All constants specify this particular formulation; their roles are explained,
but their exact values are design choices unless an explicit derivation is given.
We use $[x]_+=\max(x,0)$, $\mathbf1[A]=1$ if statement $A$ holds and zero otherwise,
and stabilizers $\varepsilon_8=10^{-8}$, $\varepsilon_9=10^{-9}$,
$\varepsilon_{12}=10^{-12}$. Vector length is $\lVert\cdot\rVert_2$; a dot product
$v\cdot w$ is the sum of products of corresponding coordinates. Matrix rows representing
items are treated as row vectors. Sums over a set of items include each item once.

\section{Start with the final score}
Define $Y(c\mid S)$ as the combined evidence, $P(c\mid S)$ as a second estimate based
mainly on product categories and text, $s_0\ge0$ as its scale, and $\kappa\ge0$ as the
strength of the second estimate. The selected settings are fixed before a new request:
\begin{equation}\label{theta}
\boxed{\theta(c\mid S)=\begin{cases}
Y(c\mid S),&\kappa=0,\\[2pt]
\displaystyle\frac{f_c}{f_c+\kappa}Y(c\mid S)
+\frac{\kappa}{f_c+\kappa}s_0P(c\mid S),&\kappa>0.
\end{cases}}
\end{equation}
\why The assumption is that more appearances make history-based evidence more dependable.
When $f_c=0$ and $\kappa>0$, the score is exactly $s_0P$; when $f_c=\kappa$, the two
estimates have equal weight. Frequency does not prove reliability for an individual product.
Section~\ref{settings} states an explicit variance assumption motivating these rational
weights. We first expand $Y$ and $P$.

\subsection{Expand $Y$: exactly what changes between its two versions}
Let $\mathcal H$ be the following ordered list of evidence sources:
\[
(\mathrm{basket,propagated,pop,cooc,household,pmi,svd,category,text,bridge}).
\]
Let $w_h\ge0$ be the selected weight of source $h$; a removed source has weight zero.
Nonnegative weights are a restriction preventing a source's contribution from being reversed.
For each source let $r_h$ be its raw score and $z_h=\mathcal Z_S[r_h]$ its standardized
score, with $\mathcal Z_S$ defined in Section~\ref{standard}. A second raw version $r_h^0$
is identical except for these two sources:
\begin{center}\small
\begin{tabular}{@{}lll@{}}\toprule
Source & $r_h^0(c\mid S)$ & $r_h(c\mid S)$ \\\midrule
basket & $p_c$ & $p_c+\delta_c$ \\
propagated & $(Lp)_c$ & $(Lp)_c+(L\delta)_c$ \\
All other sources & $r_h(c\mid S)$ & $r_h(c\mid S)$ \\\bottomrule
\end{tabular}\end{center}
Here $p_c=p(c\mid S)$ is weighted training-basket support, $\delta_c=\delta(c\mid S)$
is a sum restricted to baskets containing at least two known products, and $L$ is an
item-to-item weighting matrix. Their complete definitions follow in Section~\ref{baskets}.
Define $z_h^0=\mathcal Z_S[r_h^0]$. Let $\alpha\in[0,1]$ mix the two versions and
$\lambda_R\ge0$ weight the separately standardized pair-of-known-products score $z_R$
from Section~\ref{agreement}. Then
\begin{align}
Y_0(c\mid S)&=\sum_{h\in\mathcal H}w_hz_h^0(c\mid S),\\
Y_1(c\mid S)&=\sum_{h\in\mathcal H}w_hz_h(c\mid S),\\
Y(c\mid S)&=(1-\alpha)Y_0(c\mid S)+\alpha Y_1(c\mid S)
             +\lambda_R z_R(c\mid S).\label{Y}
\end{align}
\why Only the two explicitly identified formulas change during this mixture; the other
sources are held fixed by definition. This does not assert that other sources cannot
represent relationships involving several products. The pair term $z_R$, for example,
is controlled separately. Restricting $\alpha$ to $[0,1]$ produces an average with
nonnegative weights adding to one. These are modeling choices, not requirements of every
possible scoring function. Because the two raw columns are standardized separately,
this is not equivalent to standardizing $p+\alpha\delta$.

\subsection{Expand $P$: which sources enter and why}
Let $\mathcal A=\{h:w_h>0\}$ be the retained sources. Define the prior-source set
\begin{equation}
\mathcal P=\begin{cases}
\{\mathrm{category,text,bridge}\}\cap\mathcal A,&\text{if this intersection is nonempty},\\
\{\mathrm{category,bridge}\},&\text{otherwise}.
\end{cases}
\end{equation}
For $h\in\mathcal P$, define $\widetilde w_h=w_h$ if retained, and $\widetilde w_h=0.5$
otherwise. Thus the fallback rule may use a source whose weight in $Y$ is zero. Define
\begin{equation}\label{prior}
\boxed{P(c\mid S)=\sum_{h\in\mathcal P}\widetilde w_hz_h(c\mid S)
+\lambda_R z_R(c\mid S)+w_{\rm pop}z_{\rm pop}(c\mid S).}
\end{equation}
\why A product with no training appearances has a zero row in the training-only item
count matrices. Its category and supplied text can still distinguish it from other products.
This motivates the restricted source list. The exact list and fallback weight $0.5$ are
prescribed choices; the score equations do not establish that they are uniquely best.
Category, bridge and agreement still partly depend on training baskets, so $P$ is not
independent of $B$. Selected popularity and agreement explicitly enter both estimates.

\section{Expand standardization: what $z_h$ means}\label{standard}
For a raw source $r_h$, calculate its mean and population standard deviation across all
catalogue products for the same revealed basket $S$:
\begin{align}
\mu_h(S)&=\frac1n\sum_{i\in\I}r_h(i\mid S),\\
\sigma_h(S)&=\sqrt{\frac1n\sum_{i\in\I}(r_h(i\mid S)-\mu_h(S))^2},\\
\mathcal Z_S[r_h](c)&=\frac{r_h(c\mid S)-\mu_h(S)}{\sigma_h(S)+\varepsilon_9}.\label{z}
\end{align}
Apply this rule separately to $r_h^0$ and $r_R$ as well.
\why Counts and text similarities have different units. Equation~\eqref{z} expresses each
relative to its own catalogue average and variation. A constant column becomes zero.
The divisor here is $n$, not the sample divisor $n-1$. Known products are included in
these sums and excluded only at final ranking. These are specified conventions.
The hidden target set does not enter the calculation; neither do other test requests.
This protects standardization from label dependence, provided the raw inputs themselves
contain only information allowed at prediction time.

\section{Expand the basket and propagated scores}\label{baskets}
\subsection{The exact training-basket contribution}
Define basket size $d_e=\sum_iB_{ei}$, known-product weight $a_s=(f_s+5)^{-1/2}$,
and basket weight $v_e=(d_e+1)^{-1/2}$. Define weighted overlap $o_e$ and overlap count $k_e$:
\begin{align}
o_e(S)&=v_e\sum_{s\in S}B_{es}a_s,& k_e(S)&=\sum_{s\in S}B_{es},\\
p(c\mid S)&=\sum_e B_{ec}o_e(S),&
\delta(c\mid S)&=\sum_e B_{ec}o_e(S)\,1.5[k_e(S)-1]_+.\label{delta}
\end{align}
These give $r_{\rm basket}^0=p$ and $r_{\rm basket}=p+\delta$ in the table above.
\why $B_{ec}$ requires the candidate to occur in the contributing training basket.
The weighted overlap adds contributions from its known products, downweighting frequent
known products and large baskets. Offsets 5 and 1 keep weights finite and soften this
adjustment; their precise values and the square-root powers are choices.
For $k_e=0$ or $1$, the added contribution to $\delta$ is zero; for $k_e=2$ it is
$1.5B_{ec}o_e$; for $k_e=3$ it is $3B_{ec}o_e$. It therefore records a specific
observed overlap, not an undefined notion of group compatibility. The linear increase
and multiplier $1.5$ specify the chosen strength. With one known product, $\delta=0$.

\subsection{The exact item-to-item matrix $L$}
Let $D=\diag((f_i+5)^{-1/2})$ and $E=\diag((d_e+1)^{-1})$, where $\diag$ constructs
a diagonal matrix. Let $\off$ replace a matrix's diagonal by zero. Define
\begin{equation}\label{L}
L=\off(DB^\top EBD),\qquad
L_{ci}=\frac{\sum_e B_{ec}B_{ei}/(d_e+1)}{\sqrt{f_c+5}\sqrt{f_i+5}}\quad(c\ne i),
\qquad L_{cc}=0.
\end{equation}
Writing $p,\delta$ as column vectors over all products, the other two raw scores are
\begin{equation}
r_{\rm propagated}^0(c\mid S)=\sum_i L_{ci}p_i,\qquad
r_{\rm propagated}(c\mid S)=\sum_iL_{ci}p_i+\sum_iL_{ci}\delta_i.
\end{equation}
\why A product receives some score from other products with which it shared training
baskets. The displayed entry formula defines exactly how much. Zero diagonal excludes
immediate self-transfer. A product absent from all training baskets has a zero row in $L$.

\section{Expand household evidence}\label{household}
Let $u$ index households, and define the household-by-product count table
\begin{equation}
X_{ui}=\sum_{e:u(e)=u}B_{ei},\qquad t_u=1+\sum_iX_{ui},\qquad
\widetilde X_{ui}=X_{ui}/\sqrt{t_u}.
\end{equation}
$X_{ui}$ counts training baskets containing product $i$ for household $u$; it is not a
count of purchased units. Define $C_H=\off(\widetilde X^\top\widetilde X)$. The raw score is
\begin{equation}\label{householdscore}
\boxed{r_{\rm household}(c\mid S)=\sum_{s\in S}(C_H)_{cs}(f_s+5)^{-1/2}.}
\end{equation}
For distinct products $c,s$, $(C_H)_{cs}=\sum_u X_{uc}X_{us}/t_u$.
\why Two products can have a household connection even if they were purchased in different
baskets. The division by $t_u$ reduces the influence of households with large total counts;
the extra factor on the known product reduces frequent-product influence. These weights
are specified modeling choices. Every household contributes to the shared table $C_H$.
The current shopper's household ID does not enter equation~\eqref{householdscore}; this
score is not a personalized lookup of that shopper's history. Only training baskets
contribute to $X$.

\section{Expand popularity, pair counts, and the history summary}
Define $C=\off(B^\top B)$, $g_i=\sum_jC_{ij}+\varepsilon_9$, and
$Z_C=\sum_{i,j}C_{ij}$. For a pair observed together at least once define
\begin{equation}
M_{ij}=[\log(C_{ij}Z_C/(g_i g_j)+\varepsilon_{12})]_+;
\qquad M_{ij}=0\text{ if }C_{ij}=0.
\end{equation}
The raw sources are
\begin{equation}
r_{\rm pop}(c\mid S)=\log(1+f_c),\quad
r_{\rm cooc}(c\mid S)=\sum_{s\in S}C_{cs},\quad
r_{\rm pmi}(c\mid S)=\sum_{s\in S}M_{cs}.
\end{equation}
\why Popularity supplies a general frequency preference; $C$ measures direct shared-basket
counts. $M$ compares a pair count with its two overall shared-count masses $g_i,g_j$,
which differ from the basket frequencies $f_i,f_j$. The logarithm compresses large ratios;
the positive part discards below-expected associations rather than assigning a negative
raw penalty. This is a chosen treatment of negative association.

For the summary source, set $r=\min(128,n-2)$ and let $V_r$ contain the leading $r$ right
singular directions of $M$. For an exact decomposition these give the best rank-$r$
approximation in squared error; fitted directions may approximate them numerically.
Define row normalization $\N(F)_i=F_i/(\lVert F_i\rVert_2+\varepsilon_8)$ and set
$U=\N(MV_r)$, or $U=0$ when $M$ is zero. Then
\begin{equation}
r_{\rm svd}(c\mid S)=U_c\cdot\left(\frac1{|S|}\sum_{s\in S}U_s\right).
\end{equation}
\why The vectors summarize many pair relationships in a shorter description per product.
The dot product compares the candidate with the mean known-product summary. Width 128 is
a chosen dimension limit, not a data-derived requirement.

\section{Expand product-text evidence and its bridge}\label{text}
For each catalogue product, concatenate its department, commodity, subcommodity and brand
text. Use lowercased words and two-word phrases after removing English stop words, retaining
at most 8,000 terms. Let $n_{it}$ count term $t$ in product $i$'s text and
$\mathrm{df}_t=|\{i:n_{it}>0\}|$. Define the unnormalized text weight
\begin{equation}
A_{it}=\begin{cases}
(1+\log n_{it})\bigl[1+\log((n+1)/(\mathrm{df}_t+1))\bigr],&n_{it}>0,\\
0,&n_{it}=0.
\end{cases}
\end{equation}
Normalize each nonzero row of $A$ to length one, leaving zero rows zero, and call the
resulting sparse table $T$. Let $d_T$ be its column count, $r_T=\min(128,d_T-1)$, and
$V_T$ its leading $r_T$ right singular directions. Define the dense product-text table
$F=\N(TV_T)$. This assumes at least two retained text terms, as in this dataset.
\why The first factor in $A_{it}$ compresses repeated words; the second reduces the
relative influence of words appearing in many product descriptions. The short vectors
summarize this table. The vocabulary and dimension limits are choices. There is no
category-average subtraction in this definition of $F$.

Define the bridged feature table $F_b$ and the two selectable text-related scores by
\begin{align}
F_b&=\N\bigl(F+0.5\N(LF)\bigr),\\
r_{\rm text}(c\mid S)&=F_c\cdot\left(\frac1{|S|}\sum_{s\in S}F_s\right),\\
r_{\rm bridge}(c\mid S)&=F_{b,c}\cdot\left(\frac1{|S|}\sum_{s\in S}F_{b,s}\right).
\end{align}
\why $(LF)_c=\sum_iL_{ci}F_i$ combines descriptions of products connected to $c$ through
training baskets. The bridge mixes that normalized combination with $c$'s own description;
$0.5$ is its chosen multiplier. If $f_c=0$, then $(LF)_c=0$ but $F_c$ may still distinguish
$c$, so the bridge need not vanish. The sparse text similarity $T(T^\top x_S)$, with
$(x_S)_i=\mathbf1[i\in S]$, is not a separately weighted source in this formulation.
That exclusion is a source-menu restriction, not a statement of mathematical redundancy.

\section{Expand product-category evidence}
Let $\ell(i)$ be product $i$'s subcommodity category. Let $a,b$ index categories and define
\begin{equation}
A^{\ell}_{ea}=\sum_iB_{ei}\mathbf1[\ell(i)=a],\quad
H=\off((A^{\ell})^\top A^{\ell}),\quad
h_a=\sum_e A^{\ell}_{ea}+\varepsilon_9,\quad H_+=\sum_a h_a.
\end{equation}
For observed category pairs set
\begin{equation}
Q_{ab}=[\log(H_{ab}H_+/(h_a h_b)+\varepsilon_{12})]_+;
\qquad Q_{ab}=0\text{ if }H_{ab}=0.
\end{equation}
Then
\begin{equation}
\boxed{r_{\rm category}(c\mid S)=\sum_{s\in S}Q_{\ell(s),\ell(c)}.}
\end{equation}
\why This transfers evidence between product categories even when particular product
pairs have little history. Its marginal $h_a$ counts category appearances in training
baskets, rather than row sums of $H$. Repeated products of a category within one basket
therefore contribute repeatedly to the count table. The diagonal of $Q$ is zero: this
source uses cross-category counts, with no additional penalty for repeating a category
and no extra feature-cluster category source. These are choices specifying the model.

\section{Expand the pair-of-known-products term}\label{agreement}
Define $V=\diag((d_e+1)^{-1/2})$ and $W=\off(B^\top VBD)$, using $D$ from
Section~\ref{baskets}. Let a known product $s$ provide candidate support $q_{cs}$, and let
$A_c$ be its total support:
\begin{equation}
q_{cs}=W_{cs}+0.15[F_{b,c}\cdot F_{b,s}]_+,\qquad A_c=\sum_{s\in S}q_{cs}.
\end{equation}
For distinct known products $s,t$, define their similarity and pair weight
\begin{equation}
G_{st}=0.5\mathbf1[\ell(s)=\ell(t)]+0.5[F_{b,s}\cdot F_{b,t}]_+,
\qquad K_{st}=[1-G_{st}]_+,\qquad K_{ss}=0.
\end{equation}
The remaining raw source, standardized by equation~\eqref{z}, is
\begin{equation}\label{R}
\boxed{r_R(c\mid S)=\frac{\sum_{s\in S}\sum_{t\in S}q_{cs}K_{st}q_{ct}}{A_c+0.05}.}
\end{equation}
\why For $S=\{s,t\}$ with $s\ne t$, the numerator is exactly $2q_{cs}q_{ct}K_{st}$.
Both known products must provide nonzero support for this contribution to be positive;
$K_{st}$ reduces the contribution when they are similar. For a larger $S$, any supported
pair can contribute; matching all known products is not required. With one known product,
$K_{ss}=0$ makes the result zero. The denominator controls the scale of the product.
The constants $0.15,0.5,0.05$ set chosen feature support, similarity mixing and smoothing.
Here the descriptor is the bridged product-text vector $F_b$.

The matrix $W$ has zero diagonal, whereas $p=B^\top VBDx_S$ retains diagonal
contributions for known products. Both are evaluated over all products before the final
known-product exclusion. This distinction affects standardization and propagation even
though a known product cannot be returned as a suggestion.

\section{Expand the settings selected before prediction}\label{settings}
The weights, prior-source set, $\alpha,\lambda_R,s_0,\kappa$ are selected on held-out
basket-completion examples and then frozen. Their search grids below are finite prescribed
choices, not the only valid values or proven optimal spacings. In particular,
$\alpha\in[0,1]$ is a mathematical averaging domain, while a three-point grid is a
restricted search within that domain. Exact cutoffs, tolerances and smoothing constants
are part of this specified procedure.

\subsection{Define the comparison measure}
For held-out request $j$, let $G_j$ be its nonempty hidden-product set and $\pi_j(r)$
the product at rank $r$ after excluding $S_j$. Set $K_j=\min(20,n-|S_j|)$. Define
\begin{equation}
\mathrm{NDCG}_{20,j}=\frac{\sum_{r=1}^{K_j}\mathbf1[\pi_j(r)\in G_j]/\log_2(r+1)}
{\sum_{r=1}^{\min(20,|G_j|)}1/\log_2(r+1)},\qquad
\mathcal J=\frac1v\sum_{j=1}^v\mathrm{NDCG}_{20,j},
\end{equation}
where $v$ is the number of selection requests. This rewards hidden products near the top
of the list, with greater credit at earlier ranks. Final test labels do not enter selection.

\subsection{Weights, prior sources and the two mixture coefficients}
Starting with zero weights, visit the ordered list $\mathcal H$. At each source try
$w_h\in\{0,0.25,0.5,1,1.5\}$ with other weights fixed, skip the all-zero vector, and
accept only improvements in $\mathcal J(Y_0)$ exceeding $2\times10^{-4}$. Repeat for at
most four passes, stopping if a pass changes nothing. Record that search's incumbent score
as $J_{\rm coord}$. Try removing bridge and then category if the loss is at most
$3\times10^{-4}$, retaining at least one source. Absent weights become zero.

Next, visit $\alpha=0,0.5,1$ as the outer loop and $\lambda_R=0,0.15,0.25$ as the inner
loop. Start its score incumbent at $-\infty$ and accept a candidate when its score is at
least the incumbent minus $10^{-5}$; after acceptance replace the incumbent with that
candidate's score. A later near-tie can win, and successive accepted near-ties can move
downward slightly. This ordered rule is not simply the first global maximum.
Call the final accepted score $J_{\rm mix}$ and retain its coefficients. Construct
$\mathcal P$ as defined above, then carry $J_*=\max(J_{\rm coord},J_{\rm mix})$ into
the final comparison. This carried threshold can exceed the score of the final pruned mixture.
\why These searches determine which evidence improves the held-out ranking. The exact
visit order and acceptance tolerances define the result; they are not theoretical
statements that one combination is globally optimal.

\subsection{Prior scale and the frequency-dependent adjustment}
Using the retained settings, flatten the $Y$ and $P$ scores across all catalogue products
and selection requests. Let $\sd_0$ be population standard deviation of this flattened table.
The prior scale is fixed as
\begin{equation}
s_0=\frac{\sd_0(Y)}{\sd_0(P)+\varepsilon_9}.
\end{equation}
It can be zero. It is not recomputed from a test batch, and there is no additional
unit-scale candidate in this formulation.
\why This aligns the sizes of variations in the two estimates before combining them.

For each product with $f_i>0$, define
\begin{equation}
x_i=1/f_i,\qquad e_i=\frac1v\sum_{j=1}^v
\bigl[Y(i\mid S_j)-s_0P(i\mid S_j)\bigr]^2.
\end{equation}
Let $\bar x,\bar e$ be their means over products with positive $f_i$. Define the line's
slope $\widehat b$ and intercept $\widehat a$ by
\begin{align}
\widehat b&=\frac{\sum_i(x_i-\bar x)(e_i-\bar e)}{\sum_i(x_i-\bar x)^2},&
\widehat a&=\bar e-\widehat b\bar x,\\
\kappa_*&=\operatorname{clip}\left(\frac{\max(\widehat b,10^{-4})}{\max(\widehat a,10^{-4})},0.05,50\right).
\end{align}
Here $\operatorname{clip}(t,l,u)=\min(u,\max(l,t))$. With fewer than two products having
history, or a denominator at most $10^{-12}$, the components cannot be distinguished;
use $\kappa=0$ only. Otherwise try $\kappa\in\{0,0.05,0.25,\kappa_*,1\}$ in this
order, starting with zero and accepting only improvements exceeding $10^{-4}$ over the
carried score incumbent $J_*$. Update that incumbent after an accepted improvement.

\why The assumed relation is $\mathbb E[e_i\mid f_i]=\sigma_0^2+\sigma_e^2/f_i$.
To state the assumption precisely, let an unobserved suitability be
$\theta_i^*=s_0P_i+\xi_i$ and let $Y_i=\theta_i^*+\eta_i$, with independent zero-mean
normal errors of positive variances $\sigma_0^2$ and $\sigma_e^2/f_i$. Then, for $f_i>0$,
\begin{equation}
\mathbb E[(Y_i-s_0P_i)^2\mid f_i]=\sigma_0^2+\sigma_e^2/f_i,\qquad
\mathbb E[\theta_i^*\mid Y_i,P_i,f_i]
=\frac{f_iY_i+\kappa s_0P_i}{f_i+\kappa},\quad
\kappa=\sigma_e^2/\sigma_0^2.
\end{equation}
This motivates both the fitted line and equation~\eqref{theta}; the assumptions do not
follow from counts alone. The positivity bounds and ratio clipping produce one candidate
setting; held-out ranking comparisons choose the final value. Residuals compare the two
score estimates, not scores with binary hidden-product labels.

\section{What the training table and product inputs represent}
The preceding equations assume $B$ is fixed. In this Dunnhumby completion task it comes
from baskets of 3--20 distinct products, capped at 30,000 after a fixed shuffle. Products
appearing fewer than twice in that capped pool are removed, and baskets shorter than three
are dropped. Identical remaining product sets are deduplicated globally, retaining the
first household association. A random 70/10/20 split supplies training, selection and test
baskets; the training portion alone defines $B$ and the household table $X$.

Product categories are subcommodity labels; missing product metadata receive separate
missing-category labels. Product text and its vocabulary/short vectors use the known
catalogue. The product-frequency filter is based on the pre-split capped pool. Thus the
preparation assumes a known catalogue and does not estimate every preparation choice from
$B$ alone. The split is random, not chronological; date fields do not impose a time cutoff.
These distinctions state the information available to the formulas rather than implicitly
claiming a prediction task with unseen future products.

\section{Substitute back to obtain the recommendation}
For each candidate $c$, compute the raw terms from training baskets, training household
counts, category labels and catalogue text. Standardize them with equation~\eqref{z};
insert them into $Y$ and $P$, and then into equation~\eqref{theta}. Exclude products already
in $S$, rank the others by decreasing $\theta$, and break exact ties by increasing product
ID. For a requested positive integer $K$, return the highest $K$ candidates, or all candidates if fewer than $K$ remain. Sources
with zero weight and no role in $P$ contribute nothing. This completes the expansion
of the score without introducing an additional shopper-specific scoring term.
\end{document}
